% ============================================================================
% GLOSSARY OF KEY TERMS
% ============================================================================
\section{Glossary of Key Terms}

This section defines key terminology used throughout this document.

\subsection{Communication Protocols}

\begin{description}[style=nextline, leftmargin=2cm]
    \item[gRPC] \textbf{Google Remote Procedure Call} -- A high-performance, open-source RPC framework developed by Google. Uses HTTP/2 for transport and Protocol Buffers (Protobuf) for serialization. Enables efficient binary communication between services with built-in support for streaming.

    \item[GraphQL] A query language for APIs developed by Facebook. Unlike REST, clients specify exactly what data they need, reducing over-fetching. Supports queries (read), mutations (write), and subscriptions (real-time updates via WebSocket).

    \item[gRPC-Web] A JavaScript implementation of gRPC for browser clients. Since browsers cannot use HTTP/2 trailers directly, gRPC-Web requires a proxy (like Envoy) to translate between gRPC-Web and native gRPC.

    \item[Protobuf] \textbf{Protocol Buffers} -- Google's language-neutral, platform-neutral binary serialization format. Smaller and faster than JSON. Messages are defined in \texttt{.proto} files and compiled to language-specific code.

    \item[WebSocket] A protocol providing full-duplex communication channels over a single TCP connection. Used by GraphQL subscriptions for real-time data streaming from server to client.

    \item[REST] \textbf{Representational State Transfer} -- An architectural style for web APIs using HTTP methods (GET, POST, PUT, DELETE). Uses JSON for data exchange. Simpler but less efficient than gRPC for high-frequency communication.
\end{description}

\subsection{Middleware and Infrastructure}

\begin{description}[style=nextline, leftmargin=2cm]
    \item[DDS] \textbf{Data Distribution Service} -- A middleware protocol and API standard for data-centric publish-subscribe communication. Used by ROS 2 for inter-node communication. Provides real-time, scalable, and reliable data exchange with Quality of Service (QoS) policies.

    \item[ROS 2] \textbf{Robot Operating System 2} -- An open-source robotics middleware framework. Provides tools, libraries, and conventions for building robot applications. Uses DDS for communication between nodes.

    \item[SLAM] \textbf{Simultaneous Localization and Mapping} -- A computational technique where a robot builds a map of an unknown environment while simultaneously tracking its location within that map.

    \item[Nav2] \textbf{Navigation 2} -- The ROS 2 navigation stack. Provides autonomous navigation capabilities including path planning (A*), obstacle avoidance (DWA), and localization (AMCL).

    \item[DWA] \textbf{Dynamic Window Approach} -- A local path planning algorithm used by Nav2. Samples velocities in a ``dynamic window'' (velocities reachable within one control cycle) and scores trajectories based on proximity to goal, obstacle clearance, and velocity. Runs at 20 Hz in the STAR robot.

    \item[AMCL] \textbf{Adaptive Monte Carlo Localization} -- A probabilistic localization algorithm that uses a particle filter to track the robot's pose within a known map. Particles represent hypotheses about the robot's position; sensor observations update particle weights. SLAM Toolbox provides similar functionality during mapping.

    \item[mDNS] \textbf{Multicast DNS} -- A protocol that resolves hostnames to IP addresses within small networks without a local DNS server. Enables \texttt{star-robot.local} hostname resolution. Implemented by Avahi on Linux.

    \item[Avahi] A free, open-source implementation of mDNS/DNS-SD (Service Discovery) for Linux. Allows devices to discover each other on a local network without configuration.

    \item[Envoy] A high-performance, open-source edge and service proxy. Used in this project to translate gRPC-Web requests from browsers to native gRPC for the Spring Boot backend.
\end{description}

\subsection{Hardware Interfaces}

\begin{description}[style=nextline, leftmargin=2cm]
    \item[SPI] \textbf{Serial Peripheral Interface} -- A synchronous serial communication interface for short-distance communication. Uses 4 wires: COPI (Controller Out Peripheral In), CIPO (Controller In Peripheral Out), CLK (Clock), and CS (Chip Select). Used for RPi5-to-ESP32 communication at 10 MHz.

    \item[I2C] \textbf{Inter-Integrated Circuit} -- A two-wire serial communication protocol (SDA for data, SCL for clock). Used for connecting sensors like IMUs. Supports multiple devices on the same bus with unique addresses.

    \item[UART] \textbf{Universal Asynchronous Receiver-Transmitter} -- A serial communication protocol using TX (transmit) and RX (receive) lines. Commonly used for debugging and some sensor connections.

    \item[GPIO] \textbf{General Purpose Input/Output} -- Configurable digital pins on microcontrollers and single-board computers. Can be set as input (read sensors/buttons) or output (control LEDs/relays).

    \item[PWM] \textbf{Pulse Width Modulation} -- A technique for controlling analog devices using digital signals. The duty cycle (percentage of time the signal is HIGH) controls the effective power. Used for motor speed control.

    \item[Center-Aligned PWM] A PWM mode where the counter counts up then down, centering pulses in the period. Produces cleaner signals with reduced electromagnetic interference (EMI) compared to edge-aligned PWM. Used by the ESP32 MCPWM peripheral for motor control.

    \item[PID] \textbf{Proportional-Integral-Derivative} -- A control loop mechanism for maintaining a desired setpoint. The ESP32 runs a 1 kHz PID loop to maintain precise motor velocities based on encoder feedback.

    \item[SMBus] \textbf{System Management Bus} -- A two-wire interface based on I2C but with stricter timing requirements and additional features like packet error checking (PEC) and timeouts. Used for battery fuel gauge ICs and power management chips.

    \item[1-Wire] A serial protocol using a single data wire (plus ground) for communication. Devices are powered parasitically from the data line or via a separate power pin. Used by DS18B20 temperature sensors for thermal monitoring.
\end{description}

\subsection{Software Concepts}

\begin{description}[style=nextline, leftmargin=2cm]
    \item[PWA] \textbf{Progressive Web App} -- A web application that uses modern web technologies to deliver app-like experiences. Can work offline, be installed on devices, and receive push notifications.

    \item[Service Worker] A JavaScript worker that runs in the background, separate from the web page. Enables offline functionality by intercepting network requests and serving cached responses.

    \item[IndexedDB] A low-level browser API for storing large amounts of structured data. Used for offline telemetry storage in the PWA.

    \item[Apollo Client] A comprehensive GraphQL client for JavaScript. Provides caching, state management, and offline support for GraphQL operations.

    \item[Workbox] A set of JavaScript libraries from Google for adding offline support to web apps. Simplifies service worker creation and caching strategies.

    \item[Coroutines] A Kotlin feature for asynchronous programming. Allows writing non-blocking code in a sequential style. Used extensively in the Spring Boot gateway for handling concurrent requests.

    \item[Pi4J] A Java library for accessing GPIO, I2C, SPI, and other hardware interfaces on Raspberry Pi. Used by the gateway to communicate with the ESP32 via SPI.

    \item[FreeRTOS] A real-time operating system kernel for embedded devices. Used by the ESP32 firmware for task scheduling, ensuring the motor control loop runs at exactly 1 kHz.
\end{description}

\subsection{Protocols and Formats}

\begin{description}[style=nextline, leftmargin=2cm]
    \item[SLAMTEC Protocol] The communication protocol used by RPLiDAR sensors (like the RPLiDAR C1). A proprietary serial protocol over UART for commanding the sensor and receiving scan data, supported by the rplidar\_ros2 driver.

    \item[JSON] \textbf{JavaScript Object Notation} -- A lightweight, human-readable data interchange format. Used by GraphQL and REST APIs. Larger than binary formats like Protobuf.

    \item[HTTP/2] The second major version of HTTP. Supports multiplexing (multiple requests over one connection), header compression, and server push. Required by gRPC.

    \item[TCP/IP] \textbf{Transmission Control Protocol / Internet Protocol} -- The foundational communication protocols of the internet. Provides reliable, ordered delivery of data packets.
\end{description}

\subsection{Deployment and Operations}

\begin{description}[style=nextline, leftmargin=2cm]
    \item[systemd] The init system and service manager for modern Linux distributions. Manages service lifecycle, auto-restart on failure, and dependency ordering.

    \item[Buildroot] A tool for building embedded Linux systems. Used to create the minimal custom Linux image for the Raspberry Pi 5.

    \item[Prometheus] An open-source monitoring and alerting toolkit. Collects metrics from applications via HTTP endpoints. Used with Spring Boot Actuator for gateway monitoring.

    \item[Actuator] A Spring Boot module that exposes operational information (health, metrics, info) via HTTP endpoints. Integrates with Prometheus for metrics collection.

    \item[JVM] \textbf{Java Virtual Machine} -- The runtime environment for Java and Kotlin applications. The Spring Boot gateway runs on the JVM.
\end{description}
